% =============================================================
% บทที่ 1 - บทนำ
% =============================================================

% -----------------------------------------------------------
\section{ความเป็นมาและความสำคัญของปัญหา}
% -----------------------------------------------------------

 ข่าวสารเป็นข้อมูลสำคัญที่ช่วยให้ผู้คนรับรู้และติดตามเหตุการณ์ได้อย่างทันท่วงที โดยเฉพาะในยุคที่ข้อมูลมีการเปลี่ยนแปลงตลอดเวลา หากผู้อ่านที่ต้องการติดตามข่าวสารจากหลายแหล่งมักประสบปัญหาหลายอย่าง ได้แก่ ต้องเปิดเว็บไซต์หรือแอปพลิเคชันหลายแห่งเพื่อเปรียบเทียบข่าวเดียวกัน ทำให้เสียเวลาและเกิดความยุ่งยากในการติดตามสถานการณ์ข่าวที่เปลี่ยนรวดเร็ว นอกจากนี้ ข่าวแต่ละฉบับยังมีเนื้อหายาวและมีข้อมูลมากเกินไป ทำให้ผู้อ่านต้องใช้เวลาในการอ่านเพื่อหาประเด็นสำคัญ รวมถึงยังต้องอ่านแหล่งข่าวที่เป็นภาษาอังกฤษเพื่อให้ได้มุมมองที่หลากหลาย อีกทั้งข่าวสารที่มีปริมาณมากและหลายหมวดหมู่ยังทำให้การกรองและวิเคราะห์ข่าวเป็นเรื่องยากสำหรับผู้อ่าน โดยเฉพาะผู้ที่ต้องการติดตามข่าวสารในหมวดหมู่ที่เฉพาะเจาะจง เช่น การเมือง เศรษฐกิจ หรือเทคโนโลยี

การนำเทคโนโลยีเว็บสแครปปิง (Web Scraping) และ RSS Feed มาใช้ในการดึงข้อมูลข่าวจากเว็บไซต์ข่าวอัตโนมัติ ตลอดจนการใช้เทคนิคการประมวลผลภาษาธรรมชาติ (Natural Language Processing: NLP) ในการจำแนกหมวดหมู่ข่าวตามแต่ละประเภท และการใช้โมเดลภาษาขนาดใหญ่ (Large Language Model: LLM) ในการสรุปเนื้อหา จะช่วยลดภาระของผู้อ่านในการเปิดหลายเว็บไซต์และอ่านบทความข่าวฉบับเต็ม อีกทั้งยังสามารถนำเสนอข้อมูลสำคัญที่คัดสรรและวิเคราะห์แล้วแบบเรียลไทม์ได้ในที่เดียว

จากปัญหาดังกล่าว โครงงานนี้จึงมีแนวคิดในการพัฒนาระบบจำแนกและวิเคราะห์ข่าวสารอัตโนมัติด้วยเทคนิคการประมวลผลภาษาธรรมชาติแบบเฉพาะประเภท (Category-Specific NLP Analytics Pipeline) โดยระบบจะทำการดึงข่าวจากแหล่งข่าวหรือสำนักข่าวเป็นระยะๆแบบอัตโนมัติ จำแนกหมวดหมู่ข่าวด้วยเทคนิค NLP และโมเดลแมชชีนเลิร์นนิง สรุปเนื้อหาด้วยโมเดลภาษาขนาดใหญ่ และแสดงผลของข่าวพร้อมการวิเคราะห์บนเว็บแอปพลิเคชันแบบเรียลไทม์ผ่าน WebSocket เพื่อให้ผู้ใช้สามารถติดตามข่าวสารได้อย่างสะดวก รวดเร็ว และเหมาะสมกับความสนใจของแต่ละบุคคล

% -----------------------------------------------------------
\section{วัตถุประสงค์}
% -----------------------------------------------------------

\begin{enumerate}
    \item เพื่อศึกษาและวิเคราะห์เทคนิคการดึงข้อมูลข่าวจากเว็บไซต์ข่าวด้วยวิธีเว็บสแครปปิงและการอ่าน \\RSS Feed รวมถึงเทคนิคการประมวลผลภาษาธรรมชาติสำหรับการจำแนกหมวดหมู่ข่าวตามประเภท
    \item เพื่อออกแบบและพัฒนาระบบจำแนกและวิเคราะห์ข่าวสารอัตโนมัติด้วยเทคนิคการประมวลผลภาษาธรรมชาติแบบเฉพาะประเภท (Category-Specific NLP Analytics Pipeline) ครอบคลุมการดึงข่าว การจำแนกหมวดหมู่ การสรุปเนื้อหาด้วยโมเดลภาษาขนาดใหญ่ อัลกอริทึมแนะนำข่าว และการแสดงผลแบบเรียลไทม์
    \item เพื่อทดสอบและประเมินประสิทธิภาพของระบบในด้านความถูกต้องของการจำแนกหมวดหมู่ข่าว ความถูกต้องของการสรุปข่าว ความแม่นยำของอัลกอริทึมแนะนำข่าว และความสามารถในการแสดงผลข่าวแบบเรียลไทม์
\end{enumerate}

% -----------------------------------------------------------
\section{ขอบเขตและข้อจำกัดของระบบ}
% -----------------------------------------------------------

\subsection{ขอบเขต}

ระบบเป็นเว็บแอปพลิเคชันแบบเต็มรูปแบบ (Full-stack) ที่มีขอบเขตการทำงาน ดังนี้

\begin{enumerate}
    \item การรวบรวมข่าวอัตโนมัติ โดยระบบสามารถดึงข่าวจากแหล่งข่าวออนไลน์หลายแห่ง โดยใช้ \\เทคนิคเว็บสแครปปิง การอ่าน RSS Feed และการจำลองเบราว์เซอร์ด้วย Playwright สำหรับเว็บไซต์ที่ใช้ JavaScript
    \item การสรุปข่าวและวิเคราะห์อารมณ์ โดยระบบส่งเนื้อหาข่าวในรูปแบบ Markdown ให้ AI (LLM) ผ่าน API ของมหาวิทยาลัยขอนแก่น (KKU Gen.AI) เพื่อสร้างสรุป ประเด็นสำคัญ อารมณ์ของข่าว (Sentiment) และคำที่สำคัญ โดยสรุปเป็นภาษาเดียวกับข่าวต้นฉบับ
    \item การจำแนกหมวดหมู่ข่าว 9 หมวดหมู่หลัก ได้แก่ การเมือง (politics), เศรษฐกิจ (economy), เทคโนโลยี (tech), สุขภาพ (health), สิ่งแวดล้อม (environment), กีฬา (sports), บันเทิง (entertainment), สังคม (society) และต่างประเทศ (world) โดยใช้วิธีการตัดคำภาษาไทยด้วย PyThaiNLP ร่วมกับโมเดลการเรียนรู้ของเครื่องแบบหลายชั้น (3-Tier Cascade Classifier: Rule Cues, Calibrated LinearSVC และ WangchanBERTa)
    \item การจัดเก็บและวงจรของข้อมูล (Data Lifecycle \& Retention) ข้อมูลข่าวและผลการวิเคราะห์ทั้งหมดจะถูกจัดเก็บรวมกันต่อเนื่องนับตั้งแต่วันที่ Instance Server เริ่มทำงาน (Instance Startup Date) โดยจัดเก็บไว้ภายใน Docker Container และ Local Storage บนเครื่องแม่ข่าย AWS EC2
    \item การระบุและจัดกลุ่มประเด็นข่าวแบบออนไลน์ (Persistent Online Story Identification) ระบบสามารถตรวจจับความคล้ายคลึงเชิงความหมายแบบ Dense Vector รวมกับการยืนยันด้วยเอนทิตีชื่อที่เฉพาะเจาะจง (Named Entity Recognition) เพื่อรวบรวมข่าวจากหลายสำนักเข้าเป็นประเด็นเดียวกัน (Story Cluster) โดยกำหนดการค้นหา 72 ชั่วโมงแบบไม่กำหนดระยะเวลาหมดอายุของข่าว (Non-capping Window) ทำให้สามารถติดตามประเด็นข่าวที่ยืดเยื้อข้ามวันได้ดีขึ้น
    \item การจัดกลุ่มประเด็นข่าวหลักและประเด็นย่อย (Macro Storyline \& Sub-events) ระบบสามารถเชื่อมโยงประเด็นข่าวย่อยที่ต่อเนื่องกันเข้าสู่โครงเรื่องข่าวข่าวหลัก ในส่วนของการจำแนกประเภทกลุ่มข่าว ระบบจะประมวลผลตามหลักเสียงข้างมาก (Majority Voting) และยังเปิดให้กำหนดข้อยกเว้นเป็นกรณีพิเศษได้ทันทีเมื่อเกิดเหตุวิกฤต (Crisis Exemption)
    \item กลไกการจัดอันดับกระแสข่าว (Deterministic Multi-Signal Trending Engine) คำนวณคะแนนความนิยมจาก 6 มิติเชิงสถิติ (ปริมาณ, อัตราเร่ง, ความหลากหลายของสำนักข่าว, การลดทอนตามเวลา, ความแปลกใหม่, และพฤติกรรมผู้อ่าน) พร้อมตัวคูณลดทอนข่าวประจำวัน (Routine Penalty) โดยตัดการใช้ LLM ออกจากการคำนวณคะแนนเพื่อความรวดเร็ว
    \item มาตรการความเป็นส่วนตัว (User Privacy \& Data Flow) ระบบบันทึกเฉพาะรหัสเซสชันแบบไม่ระบุตัวตน (Anonymous Cookie ID) และพฤติกรรมการใช้งานในภาพรวม (Non-PII Aggregate Engagement เช่น การคลิกเปิดอ่าน การกดดูสรุป การบุ๊กมาร์ก) บนเบราว์เซอร์ของผู้ใช้เพื่อส่งมาคำนวณคะแนนความนิยม โดยไม่มีการจัดเก็บข้อมูลระบุตัวตนส่วนบุคคล (PII) เช่น \\ชื่อ-นามสกุล อีเมล หรือหมายเลข IP และไม่มีการส่งต่อข้อมูลพฤติกรรมผู้ใช้ไปผู้ให้บริการโมเดลภาษา (LLM) ภายนอกระบบหรือบุคคลภายนอก
    \item การแสดงผลแบบเรียลไทม์ โดยระบบจะแสดงรายการข่าวบนเว็บแอปพลิเคชันพร้อมสรุปเนื้อหา ภาพประกอบ และหมวดหมู่ และแจ้งเตือนข่าวใหม่ทันทีผ่าน WebSocket (Socket.IO) โดยไม่ต้องรีเฟรชหน้าเว็บ
    \item การอ่านข่าวในตัวเว็บ (In-app Reader Modal) โดยระบบมีหน้าต่างแสดงผลของบทความข่าวและบทวิเคราะห์ในตัวเว็บแอปพลิเคชัน เพื่อให้ผู้ใช้อ่านข่าวได้ต่อเนื่องแบบไม่ต้องเปิดแท็บใหม่
    \item การจัดเก็บข้อมูลหลายรูปแบบและสถาปัตยกรรมแบบแยกส่วน (Multi-Engine Persistence) ระบบรองรับการจัดเก็บข้อมูลในรูปแบบไฟล์ JSON, ฐานข้อมูล SQLite พร้อมโหมด \\Write-Ahead Logging (WAL) สำหรับงานอ่านเขียนประสิทธิภาพสูง และฐานข้อมูล NoSQL MongoDB ออกแบบตามหลัก SOLID/GRASP ด้วย Port/Adapter Pattern
    \item ความยืดหยุ่นในการปรับแต่งระบบ โดยระบบรองรับการเพิ่มแหล่งข่าวใหม่ได้โดยไม่ต้องแก้ไขโครงสร้างหลัก ผ่านการลงทะเบียน Source ด้วย Decorator (\texttt{@register\_source})
    \item ระบบแคชหลายระดับ (Multi-tier Cache) โดยะบบถูกออกแบบให้จัดเก็บข้อมูลข่าวและเนื้อหาไว้ในหน่วยความจำระดับ L1 (L1 Cache) (\texttt{AsyncInMemoryCache}) ร่วมกับแคชแบบมีเงื่อนไข (HTTP 304 Conditional Cache) เพื่อลดการดึงข้อมูลซ้ำและเพิ่มความรวดเร็วในการแสดงผล
    \item ระบบแดชบอร์ดและติดตามสถานะการทำงาน (Monitoring Dashboard) ระบบจัดเตรียมส่วนติดต่อผู้ใช้และ API สำหรับผู้ดูแลระบบเพื่อตรวจสอบอัตราการนำเข้าข่าว ค่าสุขภาพของแต่ละสำนักข่าว อัตราการเข้าถึงแคช และประสิทธิภาพภาพรวมของระบบ
    \item  ระบบรองรับมาตรฐาน Progressive Web App (PWA) ร่วมกับ Service Worker สำหรับการแคชหน้าเว็บเพื่อความรวดเร็ว
\end{enumerate}

\subsection{ข้อจำกัดของโครงงาน}

\begin{enumerate}
    \item ด้านทรัพยากรฮาร์ดแวร์และการติดตั้งใช้งาน ระบบถูกออกแบบและปรับแต่งให้รองรับภาระงานได้ตามขีดความสามารถของเครื่องแม่ข่าย AWS EC2 อินสแตนซ์ประเภท \texttt{c7i-flex.large} (2 vCPU, หน่วยความจำหลัก RAM ขนาด 4 GiB และ Swap Buffer ขนาด 4 GB) ซึ่งอาจมีข้อจำกัดในการรองรับผู้ใช้งานพร้อมกันจำนวนมาก (Concurrent Users) หรือการรันเบราว์เซอร์ Playwright จำนวนหลายอินสแตนซ์พร้อมกัน
    \item ด้านชุดข้อมูลสำหรับฝึกโมเดล (Training Dataset Constraints) โมเดลการเรียนรู้ของเครื่องสำหรับการจำแนกหมวดหมู่ข่าวได้รับการฝึกสอนและประเมินผลบนชุดข้อมูลจำกัดจากคลังข้อมูล ThaiSum และ Prachathai-67k ซึ่งอาจมีความครอบคลุมของคำศัพท์เฉพาะกลุ่มหรือโครงสร้างภาษาตามลักษณะของแหล่งข่าวดังกล่าว
    \item ด้านความหลากหลายของการพาดหัวข่าวและแท็กหมวดหมู่ข้ามสำนักพิมพ์ ข่าวเหตุการณ์เดียวกันจากต่างสำนักข่าวอาจมีการพาดหัวคนละมุมมอง หรือมีโครงสร้าง URL และ Tag หมวดหมู่ที่แตกต่างกัน ซึ่งระบบแก้ไขด้วยการใช้ฟีเจอร์แบบไฮบริด (หัวข้อ, เนื้อหาย่อ, คำสำคัญ, \\เอนทิตี) และการหาฉันทามติของหมวดหมู่ในระดับกลุ่มข่าว (Cluster-Level Category Consensus)
    \item ระบบรองรับการสรุปและแสดงผลข่าวในภาษาไทยและภาษาอังกฤษเท่านั้น ตามภาษาเดิมของเนื้อหาข่าวแต่ละฉบับ
    \item โครงสร้าง HTML หรือ RSS Feed ของเว็บไซต์แหล่งข่าวอาจเปลี่ยนแปลงได้ ทำให้ระบบต้องได้รับการปรับปรุงให้สอดคล้อง
    \item การสรุปข่าวต้องอาศัยการเชื่อมต่ออินเทอร์เน็ตและ API ของโมเดล LLM (KKU Gen.AI) หากบริการขัดข้องหรือหมดโควตาการใช้งาน ระบบจะทำงานในโหมดข้อความสรุปตั้งต้นแทน
    \item รอบแรกของการดึงข่าวอาจใช้เวลา 2--5 นาที เนื่องจากต้องดึงเนื้อหาบทความของแต่ละข่าวเพื่อใช้ในการสรุป
\end{enumerate}

% -----------------------------------------------------------
\section{ประโยชน์ที่คาดว่าจะได้รับ}
% -----------------------------------------------------------

\begin{enumerate}
    \item ได้ระบบเว็บแอปพลิเคชันจำแนกและวิเคราะห์ข่าวสารอัตโนมัติแบบเรียลไทม์ ที่ช่วยให้ผู้ใช้ติดตามข่าวสารจากหลายแหล่งได้ในที่เดียว รวดเร็ว และเหมาะสมกับความสนใจของแต่ละบุคคล
    \item ได้แนวทางในการประยุกต์ใช้เทคนิคเว็บสแครปปิง การอ่าน RSS Feed โมเดลภาษาขนาดใหญ่ \\อัลกอริทึมการกำหนดคุกกี้เพื่อการปรับแต่งเฉพาะบุคคล และอัลกอริทึมแนะนำข่าว เพื่อสร้างระบบวิเคราะห์ข้อมูลข่าวอัตโนมัติ ซึ่งสามารถนำไปพัฒนาต่อยอดกับข้อมูลประเภทอื่น
    \item ได้องค์ความรู้ด้านการประมวลผลภาษาธรรมชาติ การจำแนกข้อความด้วยกฎคำสำคัญร่วมกับโมเดลแมชชีนเลิร์นนิง เทคนิคการสรุปข้อความด้วยโมเดลภาษาขนาดใหญ่ และการพัฒนาระบบเว็บแบบเรียลไทม์
\end{enumerate}
